%%%%%%
%
% $Autor: Srikanth $
% $Datum: 2020-01-18 11:15:45Z $
% $Pfad: githubtemplate/Template/Presentations/Template/slides/rename.tex $
% $Version: 4620 $
%
%
% !TeX encoding = utf8
% !TeX root = Rename
%
%%%%%%

%%%%%%%%%%%%%%%%%%%%%%%%%%%%%%%%%%%%%%%%%%%%%%%%%%%%%%%%%%%%
\Mysection{Feature-Based Ground Texture Localization}

\STANDARD{Idea}{%
Ground-texture localization relies on:
\begin{itemize}
    \item Detecting stable keypoints on asphalt, concrete, tiles \cite{Agrawal:2008,Bay:2006,Alcantarilla:2013}.
    \item Extracting descriptors such as SIFT, SURF, AKAZE, BRIEF, ORB \cite{Lowe:2004,Calonder:2010,Rublee:2011}.
    \item Matching descriptors against a map or between frames \cite{Matas:2004}.
    \item Pose estimation via RANSAC \cite{Desouza:2002}.
\end{itemize}
}

\STANDARD{Hyperparameters}{%
\begin{itemize}
    \item Number of keypoints (200–2000) \cite{Schmid:2019}.
    \item Descriptor size: SIFT (128), ORB (32) \cite{Lowe:2004,Rublee:2011}.
    \item Matching threshold: Lowe ratio 0.7–0.9 \cite{Lowe:2004}.
    \item RANSAC iterations (100–1000) \cite{Desouza:2002}.
\end{itemize}
}

\STANDARD{Library}{%
\begin{itemize}
    \item OpenCV: SIFT, SURF, ORB, BRISK, AKAZE \cite{Bradski:2000}.
    \item libv4l + OpenCV for downward camera capture.
\end{itemize}
}

\STANDARD{Applications}{%
\begin{itemize}
    \item Indoor AMRs (autonomous mobile robots) \cite{Kelly:2007}.
    \item Industrial AGVs without infrastructure \cite{Zhang:2019}.
    \item Warehouse floor navigation \cite{Chen:2018}.
\end{itemize}


%%% TikZ FIGURE 1
\begin{figure}[h]
\centering
\begin{tikzpicture}[scale=0.5, node distance=0.5cm,
proc/.style={rectangle, draw, rounded corners, minimum width=1.5cm, minimum height=0.9cm}]
\node[proc] (img) {Image};
\node[proc, right=of img] (kp) {Keypoints};
\node[proc, right=of kp] (desc) {Descriptors};
\node[proc, right=of desc] (match) {Matching};
\node[proc, right=of match] (pose) {Pose};

\draw[->] (img) -- (kp);
\draw[->] (kp) -- (desc);
\draw[->] (desc) -- (match);
\draw[->] (match) -- (pose);
\end{tikzpicture}
\caption{Feature-based localization pipeline}
\end{figure}
}
%%%%%%%%%%%%%%%%%%%%%%%%%%%%%%%%%%%%%%%%%%%%%%%%%%%%%%%%%%%%
\Mysection{Appearance-Based Localization}

\STANDARD{Idea}{%
\begin{itemize}
    \item Uses correlation or template matching between patches \cite{Kelly:2007,Zaman:2007}.
    \item Computes pixel displacement using NCC/SSD \cite{Pratt:2007}.
    \item Converts displacement to metric motion with planar assumption.
\end{itemize}
}

\STANDARD{Hyperparameters}{%
\begin{itemize}
    \item Patch size (16–64 px) \cite{Aqel:2016}.
    \item Search window radius (10–50 px).
    \item NCC/SSD threshold for valid match.
\end{itemize}
}

\STANDARD{Library}{%
\begin{itemize}
    \item OpenCV: matchTemplate(), NCC/SSD.
\end{itemize}
}

\STANDARD{Applications}{%
\begin{itemize}
    \item Incremental localization for factory robots \cite{Kelly:2007}.
    \item Floor tracking for simple AGVs \cite{Zaman:2007}.
\end{itemize}


%%% TikZ FIGURE 2
\begin{figure}[h]
\centering
\begin{tikzpicture}[scale=0.5, node distance=0.5cm,
proc/.style={rectangle, draw, rounded corners, minimum width=1.5cm, minimum height=0.9cm}]
\node[proc] (ref) {Reference Patch};
\node[proc, right=of ref] (search) {Search Window};
\node[proc, right=of search] (corr) {Correlation};
\node[proc, right=of corr] (motion) {Motion Estimate};

\draw[->] (ref) -- (search);
\draw[->] (search) -- (corr);
\draw[->] (corr) -- (motion);
\end{tikzpicture}
\caption{Appearance-based localization flow }
\end{figure}
}
%%%%%%%%%%%%%%%%%%%%%%%%%%%%%%%%%%%%%%%%%%%%%%%%%%%%%%%%%%%%
\Mysection{Feature-Based Visual Odometry}

\STANDARD{Idea}{%
\begin{itemize}
    \item Detect keypoints in consecutive frames \cite{Agrawal:2008,Alcantarilla:2013}.
    \item Match descriptors (ORB, BRISK, AKAZE) \cite{Rublee:2011,Leutenegger:2011}.
    \item Estimate frame-to-frame motion with RANSAC \cite{Desouza:2002}.
\end{itemize}
}

\STANDARD{Hyperparameters}{%
\begin{itemize}
    \item Keypoint budget per frame (500–1500).
    \item Hamming distance threshold for binary descriptors.
    \item RANSAC inlier threshold (1–3 px).
\end{itemize}
}

\STANDARD{Library}{%
\begin{itemize}
    \item OpenCV: ORB-SLAM2 style pipeline \cite{MurArtal:2017}.
\end{itemize}
}

\STANDARD{Applications}{%
\begin{itemize}
    \item Short-range floor-based odometry \cite{Nakashima:2019,Kozak:2016}.
\end{itemize}


%%% FIGURE 3
\begin{figure}[h]
\centering
\begin{tikzpicture}[scale=0.85, node distance=1cm,
proc/.style={rectangle, draw, rounded corners, minimum width=1.5cm, minimum height=0.9cm}]
\node[proc] (f1) {Frame k};
\node[proc, below=1cm of f1] (f2) {Frame k+1};
\node[proc, right=1cm of f1] (match) {Matching};
\node[proc, right=of match] (ransac) {RANSAC};
\node[proc, right=of ransac] (upd) {Odometry};

\draw[->] (f1) -- (match);
\draw[->] (f2) -- (match);
\draw[->] (match) -- (ransac);
\draw[->] (ransac) -- (upd);
\end{tikzpicture}
\caption{Feature-based visual odometry (scale=0.85)}
\end{figure}
}

%%%%%%%%%%%%%%%%%%%%%%%%%%%%%%%%%%%%%%%%%%%%%%%%%%%%%%%%%%%%
\Mysection{Map-Based Absolute Localization}

\STANDARD{Idea}{%
\begin{itemize}
    \item Offline mapping of ground features \cite{Zhang:2019}.
    \item Online matching of current image to map keypoints \cite{Chen:2018}.
    \item RANSAC pose recovery \cite{Schmid:2019}.
\end{itemize}
}

\STANDARD{Hyperparameters}{%
\begin{itemize}
    \item Map feature density (5–20 keypoints/m\(^2\)).
    \item Descriptor type for map (SIFT/AKAZE) \cite{Zhang:2019}.
\end{itemize}
}

\STANDARD{Library}{%
\begin{itemize}
    \item OpenCV + FLANN \cite{Bradski:2000}.
\end{itemize}
}

\STANDARD{Applications}{%
\begin{itemize}
    \item Autonomous warehouse robots \cite{Zhang:2019}.
    \item Large-area indoor localization \cite{Cornick:2016}.
\end{itemize}

%%% FIGURE 4
\begin{figure}[h]
\centering
\begin{tikzpicture}[scale=0.9, node distance=0.5cm,
proc/.style={rectangle, draw, rounded corners, minimum width=1.5cm, minimum height=0.9cm},
db/.style={cylinder, draw, aspect=0.3, minimum height=1cm, minimum width=0.5cm}]
\node[proc] (map) {Offline Map};
\node[db, right=of map] (dbx) {DB};
\node[proc, right=of dbx] (query) {Query Frame};
\node[proc, right=of query] (retr) {Retrieval};
\node[proc, right=of retr] (pose) {Pose};

\draw[->] (map) -- (dbx);
\draw[->] (query) -- (retr);
\draw[->] (dbx) -- (query);
\draw[->] (retr) -- (pose);
\end{tikzpicture}
\caption{Map-based absolute localization}
\end{figure}
}
%%%%%%%%%%%%%%%%%%%%%%%%%%%%%%%%%%%%%%%%%%%%%%%%%%%%%%%%%%%%
\Mysection{Example Code for Ground Texture Algorithms}

\STANDARD{Full File}{
\lstinputlisting[
  language=Python,
  style=mystyle,
  caption={Full Ground Texture Algorithm Module (Ground.py)}
]{Code/Ground.py}
}

\STANDARD{Next Steps}
{
	\begin{itemize}
		\item Explore each algorithm and find out its hyperparameters and metrics.
		\item Find more research paper based on each algorithm. 
		\item If the JetBot system is facing hardware constraints then find other methods to implement a ground localization algorithm.
	\end{itemize}
}
